\scanpage{70}
\mspage{- 7 -}
\margin{n° 384}
sur $T^\circ$ et est strictement positive pour $f\ne 0$. Par suite, $v$ est enveloppe supérieure de fonctions continues et est donc semi-continue inférieurement.

Soit maintenant $g$ une fonction numérique positive, $\mu$-intégrable sur $G$. D'après la prop. 5 du ch. VII, § 2, la fonction $g^{b(\alpha)}$ sur $Q_\alpha$ définie par
\[
g^{b(\alpha)}\circ\varpi_\alpha(x)=\int^*g(xs)\,d\alpha(s)
\]
est finie presque partout (pour $\mu/\alpha$) et est $\mu/\alpha$-intégrable ; de plus, on a $\mu(g)=(\mu/\alpha)(g^{b(\alpha)})$.

PROPOSITION 5.- \uline{Soit} $g$ \uline{une fonction numérique positive} $\mu$-\uline{intégrable sur} $G$ \uline{et soit} $T^\circ(g)$ \uline{l'ensemble des} $\alpha\in T^\circ$ \uline{telles que} $\inf g^{b(\alpha)}\geq 1$. \uline{La restriction de} $v$ \uline{à} $T^\circ(g)$ \uline{est continue}.

D'après la prop. 4, il suffit de montrer que la restriction de $v$ à $T^\circ(g)$ est semi-continue supérieurement. Soit donc $\unsure{\Phi|\mathfrak F}$ un filtre sur $T^\circ(g)$, convergeant vers un point $\alpha\in T^\circ(g)$. Soit $\varepsilon$ un nombre réel $>0$ et soit $h\in\mathcal K_+(G)$ avec $\mu(|g-h|)<\varepsilon$. Soit $K$ le support de $h$ et soit $F$ une fonction de $\mathcal K(Q_\alpha)$, comprise entre 0 et 1, avec $F=1$ sur le compact $\varpi_\alpha(K)$. Il existe une fonction $f\in\mathcal K_+(G)$ avec $f^{b(\alpha)}=F$ (ch. VII, § 2, n° 1, prop. 2) : on a donc $f*\alpha\leq 1$ partout et $f*\alpha=1$ sur $K$. De plus, pour tout $\beta\in T^\circ$, on a :
\[
(10)\qquad v(\alpha)\geq\mu/\alpha(F)=\mu/\alpha(f^{b(\alpha)})=\mu(f)=(\mu/\beta)(f^{b(\beta)})
\]
et $f^{b(\beta)}\circ\varpi_\beta=f*\beta$. Les propriétés de continuité de l'application $\beta\longrightarrow f*\beta$ montrent qu'il existe un ensemble $A\in\unsure{\Phi|\mathfrak F}$ tel que $f*\beta(x)>1-\varepsilon$ pour tout $x\in K$ et $\beta\in A$. Si $\beta\in A$, on a donc $f^{b(\beta)}>1-\varepsilon$ sur $\varpi_\beta(K)$ et par suite :
\[
(11)\qquad (\mu/\beta)(f^{b(\beta)})>(1-\varepsilon)(\mu/\beta)(\varpi_\beta(K)).
\]
Posons $X_\beta=Q_\beta-\varpi_\beta(K)$ et $Y_\beta=\varpi^{-1}(X_\beta)$ : on a $K\subset G-Y_\beta$ et $h=0$ sur $Y_\beta$.
Soit $\varphi_\beta$ (resp. $\psi_\beta$) la fonction caractéristique de l'ouvert $X_\beta$ (resp. $Y_\beta$) ; on a :
\margin{69}
